\documentclass{beamer}
\usetheme{Madrid}
\usecolortheme{dolphin}
\usefonttheme{structurebold}
\logo{\rule{0.6cm}{0.6cm}}
\title{Madrid with Dolphin}
\author{C. Coder}
\institute{FlashTeX}
\date{October 2026}
\begin{document}
\begin{frame}\titlepage\end{frame}
\section{Theorems}
\begin{frame}{Definitions and theorems}
  \begin{definition}
    A \emph{prime} is a natural number greater than one with no divisors
    other than one and itself.
  \end{definition}
  \begin{theorem}[Euclid]
    There are infinitely many primes.
  \end{theorem}
  \begin{proof}
    Suppose finitely many, multiply them and add one.
  \end{proof}
\end{frame}
\begin{frame}<1-2>[label=steps]{Steps, revisited later}
  \begin{enumerate}
    \item<1-> Collect the primes.
    \item<2-> Multiply them.
    \item<3-> Add one.
  \end{enumerate}
\end{frame}
\section{Examples}
\begin{frame}{Examples}
  \begin{example}
    $2, 3, 5, 7, 11, 13$.
  \end{example}
  \begin{corollary}
    There is no largest prime.
  \end{corollary}
\end{frame}
\againframe<3>{steps}
\end{document}
